\documentclass[a4paper,12pt]{article}
\usepackage[utf8]{inputenc}
\usepackage[spanish]{babel}
\usepackage{amsmath}
\usepackage{amsfonts}
\usepackage{amssymb}
\usepackage{siunitx}
\usepackage{geometry}
\usepackage{enumitem}

\geometry{margin=2.5cm}

\title{Examen de Física I - Forma B}
\author{Nombre: \underline{\hspace{8cm}} \quad Matrícula: \underline{\hspace{3cm}}}
\date{Fecha: \underline{\hspace{4cm}}}

\begin{document}

\maketitle

\section{Parte I: Teoría (4 puntos - 1 pt c/u)}
Responda brevemente a las siguientes preguntas:

\begin{enumerate}[itemsep=2em]
    \item \textbf{Etimología}: ¿De qué vocablo griego proviene la palabra ``Física'' y qué significa?
    \item \textbf{Redondeo}: ¿Qué criterio define la cantidad de decimales en el resultado de una suma o resta?
    \item \textbf{Notación Científica}: Describa la principal utilidad de este sistema de representación.
    \item \textbf{Magnitudes Derivadas}: Explique qué son las magnitudes derivadas y proporcione dos ejemplos incluyendo sus unidades.
\end{enumerate}

\vspace{2cm}

\section{Parte II: Práctica (6 puntos - 1.5 pts c/u)}
Resuelva los siguientes problemas mostrando todo su procedimiento. Redondee según las reglas de cifras significativas aprendidas.

\begin{enumerate}[itemsep=4em]
    \item \textbf{Conversión de Caudal}: Convierta $45,000 \text{ cm}^3/\text{min}$ a galones por segundo (gal/s). (Use $1\text{ gal} = 3785.41\text{ cm}^3$).
    \item \textbf{Conversión de Velocidad}: Convierta \SI{2500}{\km\per\second} a millas por minuto (mi/min). (Use $1\text{ mi} = 1.60934\text{ km}$).
    \item \textbf{Notación Científica}: Realice la siguiente división y exprese el resultado en notación científica: $(8.4 \times 10^{15}) / (2.1 \times 10^6)$.
    \item \textbf{Cifras Significativas}: Realice la siguiente resta y redondee correctamente: $450.85 - 12.5$.
\end{enumerate}

\end{document}
